\chapter{PCA und Dimensionsreduktion}

\section{Ziel dieses Kapitels}

In vielen Anwendungen des maschinellen Lernens sind Daten hochdimensional.
Ein Bild mit \(28 \times 28\) Pixeln kann zum Beispiel als Vektor in

\[
    \mathbb{R}^{784}
\]

dargestellt werden.
Moderne Datensätze können sogar Tausende oder Millionen Merkmale enthalten.

Hohe Dimensionen sind aber nicht immer hilfreich.
Oft liegen die Daten näherungsweise auf einer niedrigdimensionalen Struktur.
Dimensionsreduktion (dimensionality reduction) versucht, eine niedrigdimensionale Darstellung der Daten zu finden, die möglichst viel relevante Information erhält.

\begin{conceptbox}
Dimensionsreduktion (dimensionality reduction) bedeutet, Daten aus einem hochdimensionalen Raum in einen niedrigdimensionalen Raum abzubilden.

Dabei soll möglichst viel wichtige Struktur der Daten erhalten bleiben.
\end{conceptbox}

In diesem Kapitel betrachten wir vor allem Hauptkomponentenanalyse (principal component analysis), kurz PCA.
PCA ist eine der wichtigsten linearen Methoden zur Dimensionsreduktion.

Nach diesem Kapitel solltest du folgende Fragen beantworten können:

\begin{itemize}
    \item Warum braucht man Dimensionsreduktion (dimensionality reduction)?
    \item Was ist der Fluch der Dimensionalität (curse of dimensionality)?
    \item Was ist PCA (principal component analysis)?
    \item Was bedeutet Zentrierung der Daten?
    \item Was ist eine Kovarianzmatrix (covariance matrix)?
    \item Was sind Eigenwerte und Eigenvektoren?
    \item Wie findet PCA die Richtungen maximaler Varianz?
    \item Wie projiziert man Daten auf Hauptkomponenten?
    \item Wie rekonstruiert man Daten aus einer reduzierten Darstellung?
    \item Was ist Rekonstruktionsfehler (reconstruction error)?
    \item Wie wählt man die Anzahl der Hauptkomponenten?
    \item Was ist erklärte Varianz (explained variance)?
    \item Wie hängt PCA mit Singular Value Decomposition zusammen?
\end{itemize}

\section{Motivation für Dimensionsreduktion}

Ein Datenpunkt sei ein Vektor

\[
    \vec{x}
    =
    \begin{pmatrix}
        x_1 \\
        x_2 \\
        \vdots \\
        x_D
    \end{pmatrix}
    \in \mathbb{R}^D.
\]

Wenn \(D\) groß ist, sprechen wir von hochdimensionalen Daten.

Beispiele:

\begin{itemize}
    \item Bilder mit vielen Pixeln,
    \item Genexpressionsdaten mit vielen Genen,
    \item medizinische Zeitreihen mit vielen Messwerten,
    \item Textdaten mit vielen Wortmerkmalen,
    \item Sensordaten mit vielen Kanälen.
\end{itemize}

Dimensionsreduktion versucht, eine Abbildung zu finden:

\[
    \vec{x}
    \in
    \mathbb{R}^D
    \quad
    \longrightarrow
    \quad
    \vec{z}
    \in
    \mathbb{R}^M,
\]

wobei

\[
    M < D.
\]

Der Vektor \(\vec{z}\) ist die niedrigdimensionale Darstellung (low-dimensional representation).

\begin{definitionbox}
Dimensionsreduktion (dimensionality reduction) sucht eine Darstellung

\[
    \vec{z}\in\mathbb{R}^M
\]

für Datenpunkte

\[
    \vec{x}\in\mathbb{R}^D,
\]

wobei

\[
    M<D.
\]
\end{definitionbox}

\section{Warum hohe Dimensionen problematisch sind}

Hohe Dimensionen können mehrere Probleme verursachen.

\subsection{Rechenaufwand}

Viele Algorithmen werden mit steigender Dimension teurer.
Matrixoperationen, Distanzberechnungen und Optimierung können deutlich langsamer werden.

\subsection{Speicherbedarf}

Ein Datensatz mit \(N\) Datenpunkten und \(D\) Merkmalen kann als Matrix geschrieben werden:

\[
    X \in \mathbb{R}^{N\times D}.
\]

Wenn \(N\) und \(D\) groß sind, braucht diese Matrix viel Speicher.

\subsection{Rauschen und irrelevante Merkmale}

Nicht alle Merkmale sind gleich wichtig.
Manche Merkmale enthalten vor allem Rauschen (noise) oder redundante Information.

\subsection{Visualisierung}

Menschen können Daten in zwei oder drei Dimensionen gut visualisieren.
Daten in

\[
    \mathbb{R}^{100}
\]

oder

\[
    \mathbb{R}^{1000}
\]

sind dagegen schwer direkt darzustellen.

\begin{conceptbox}
Dimensionsreduktion hilft bei:

\begin{itemize}
    \item Kompression,
    \item Visualisierung,
    \item Rauschreduktion,
    \item schnellerer Berechnung,
    \item besserer Generalisierung,
    \item Erkennen wichtiger Datenrichtungen.
\end{itemize}
\end{conceptbox}

\section{Der Fluch der Dimensionalität}

Der Fluch der Dimensionalität (curse of dimensionality) beschreibt Probleme, die auftreten, wenn die Dimension sehr groß wird.

Eine wichtige Intuition:
In hohen Dimensionen wird der Raum extrem groß.
Bei gleicher Anzahl von Datenpunkten wird der Raum sehr dünn besetzt.

\begin{examplebox}
Betrachte das Intervall \([0,1]\) in einer Dimension.
Wenn wir \(10\) Punkte gleichmäßig verteilen, ist der Abstand ungefähr klein.

In zwei Dimensionen hat das Einheitsquadrat \([0,1]^2\).
Um in jeder Richtung \(10\) Positionen zu haben, brauchen wir

\[
    10^2 = 100
\]

Punkte.

In \(D\) Dimensionen brauchen wir

\[
    10^D
\]

Punkte.
\end{examplebox}

Für \(D=10\) wären das bereits

\[
    10^{10}
\]

Punkte.

\begin{conceptbox}
Der Fluch der Dimensionalität (curse of dimensionality) bedeutet:
Mit steigender Dimension wächst der Raum so schnell, dass endliche Datensätze sehr dünn werden.

Dadurch werden Distanzen, Dichteschätzungen und Nachbarschaftsbegriffe schwieriger.
\end{conceptbox}

\section{Lineare Dimensionsreduktion}

PCA ist eine lineare Methode.
Das bedeutet:
Die niedrigdimensionale Darstellung entsteht durch Projektion auf lineare Richtungen.

Wir suchen \(M\) Richtungen

\[
    \vec{u}_1,\dots,\vec{u}_M
    \in
    \mathbb{R}^D,
\]

und projizieren einen Datenpunkt \(\vec{x}\) auf diese Richtungen.

Für eine Richtung \(\vec{u}_j\) ist die Projektion:

\[
    z_j
    =
    \vec{u}_j^{\top}\vec{x}.
\]

Die reduzierte Darstellung ist:

\[
    \vec{z}
    =
    \begin{pmatrix}
        z_1 \\
        \vdots \\
        z_M
    \end{pmatrix}.
\]

Wenn wir die Richtungen als Spalten einer Matrix sammeln,

\[
    U_M
    =
    \begin{pmatrix}
        \vec{u}_1 & \vec{u}_2 & \dots & \vec{u}_M
    \end{pmatrix}
    \in
    \mathbb{R}^{D\times M},
\]

dann ist:

\[
    \vec{z}
    =
    U_M^{\top}\vec{x}.
\]

\begin{definitionbox}
Eine lineare Dimensionsreduktion projiziert Daten mit einer Matrix \(U_M\) auf einen niedrigdimensionalen Raum:

\[
    \vec{z}
    =
    U_M^{\top}\vec{x},
\]

wobei

\[
    U_M\in\mathbb{R}^{D\times M}
\]

und

\[
    M<D.
\]
\end{definitionbox}

\section{Zentrierung der Daten}

PCA wird auf zentrierten Daten durchgeführt.
Das bedeutet:
Man zieht von jedem Datenpunkt den empirischen Mittelwert ab.

Gegeben seien Datenpunkte

\[
    \vec{x}_1,\dots,\vec{x}_N
    \in
    \mathbb{R}^D.
\]

Der empirische Mittelwert ist:

\[
    \vec{\mu}
    =
    \frac{1}{N}
    \sum_{n=1}^{N}
    \vec{x}_n.
\]

Die zentrierten Daten sind:

\[
    \tilde{\vec{x}}_n
    =
    \vec{x}_n-\vec{\mu}.
\]

\begin{definitionbox}
Zentrierung (centering) bedeutet, den empirischen Mittelwert von allen Datenpunkten abzuziehen:

\[
    \tilde{\vec{x}}_n
    =
    \vec{x}_n-\vec{\mu}.
\]
\end{definitionbox}

Nach der Zentrierung gilt:

\[
    \frac{1}{N}
    \sum_{n=1}^{N}
    \tilde{\vec{x}}_n
    =
    \vec{0}.
\]

\begin{notebox}
PCA sucht Richtungen der Varianz um den Mittelwert.
Deshalb müssen die Daten vorher zentriert werden.

Wenn man nicht zentriert, kann die erste Hauptkomponente stark durch die Lage des Mittelwerts beeinflusst werden.
\end{notebox}

\section{Datenmatrix}

Wir sammeln die zentrierten Daten in einer Matrix:

\[
    \tilde{X}
    =
    \begin{pmatrix}
        \tilde{\vec{x}}_1^{\top} \\
        \tilde{\vec{x}}_2^{\top} \\
        \vdots \\
        \tilde{\vec{x}}_N^{\top}
    \end{pmatrix}
    \in
    \mathbb{R}^{N\times D}.
\]

Jede Zeile ist ein Datenpunkt.
Jede Spalte ist ein Merkmal.

\begin{notebox}
Je nach Literatur liegen Datenpunkte manchmal als Spalten statt als Zeilen in der Datenmatrix.
Dann ändern sich manche Matrixformeln durch Transposition.

In diesem Kapitel verwenden wir:

\[
    \tilde{X}\in\mathbb{R}^{N\times D}.
\]
\end{notebox}

\section{Kovarianzmatrix}

Die Kovarianzmatrix beschreibt, wie Merkmale gemeinsam variieren.
Für zentrierte Daten ist die empirische Kovarianzmatrix:

\[
    S
    =
    \frac{1}{N}
    \tilde{X}^{\top}\tilde{X}.
\]

Damit gilt:

\[
    S \in \mathbb{R}^{D\times D}.
\]

\begin{definitionbox}
Die empirische Kovarianzmatrix (empirical covariance matrix) zentrierter Daten ist:

\[
    S
    =
    \frac{1}{N}
    \sum_{n=1}^{N}
    \tilde{\vec{x}}_n\tilde{\vec{x}}_n^{\top}.
\]

In Matrixform:

\[
    S
    =
    \frac{1}{N}
    \tilde{X}^{\top}\tilde{X}.
\]
\end{definitionbox}

Der Eintrag \(S_{ij}\) ist die empirische Kovarianz zwischen Merkmal \(i\) und Merkmal \(j\).

\[
    S_{ij}
    =
    \frac{1}{N}
    \sum_{n=1}^{N}
    \tilde{x}_{ni}\tilde{x}_{nj}.
\]

\begin{conceptbox}
Die Diagonale der Kovarianzmatrix enthält Varianzen.

Die Nicht-Diagonaleinträge enthalten Kovarianzen zwischen verschiedenen Merkmalen.
\end{conceptbox}

\section{Symmetrie und positive Semidefinitheit}

Die Kovarianzmatrix \(S\) ist symmetrisch:

\[
    S^{\top}
    =
    S.
\]

Außerdem ist sie positiv semidefinit.
Das bedeutet:

\[
    \vec{v}^{\top}S\vec{v}
    \geq
    0
\]

für alle

\[
    \vec{v}\in\mathbb{R}^D.
\]

Warum?

\[
    \vec{v}^{\top}S\vec{v}
    =
    \vec{v}^{\top}
    \left(
        \frac{1}{N}
        \tilde{X}^{\top}\tilde{X}
    \right)
    \vec{v}.
\]

Also:

\[
    \vec{v}^{\top}S\vec{v}
    =
    \frac{1}{N}
    \vec{v}^{\top}
    \tilde{X}^{\top}\tilde{X}
    \vec{v}.
\]

Das ist:

\[
    \vec{v}^{\top}S\vec{v}
    =
    \frac{1}{N}
    (\tilde{X}\vec{v})^{\top}
    (\tilde{X}\vec{v}).
\]

Damit:

\[
    \vec{v}^{\top}S\vec{v}
    =
    \frac{1}{N}
    \|\tilde{X}\vec{v}\|^2
    \geq
    0.
\]

\begin{definitionbox}
Eine Matrix \(S\) heißt positiv semidefinit (positive semidefinite), wenn gilt:

\[
    \vec{v}^{\top}S\vec{v}\geq 0
\]

für alle Vektoren \(\vec{v}\).
\end{definitionbox}

\section{Varianz in einer Richtung}

PCA sucht Richtungen, in denen die Daten möglichst stark variieren.
Sei \(\vec{u}\) eine Richtung mit Norm \(1\):

\[
    \|\vec{u}\|=1.
\]

Die Projektion eines zentrierten Datenpunkts auf \(\vec{u}\) ist:

\[
    z_n
    =
    \vec{u}^{\top}\tilde{\vec{x}}_n.
\]

Die empirische Varianz der projizierten Daten ist:

\[
    \frac{1}{N}
    \sum_{n=1}^{N}
    z_n^2.
\]

Einsetzen von \(z_n\):

\[
    \frac{1}{N}
    \sum_{n=1}^{N}
    (\vec{u}^{\top}\tilde{\vec{x}}_n)^2.
\]

Da

\[
    (\vec{u}^{\top}\tilde{\vec{x}}_n)^2
    =
    \vec{u}^{\top}
    \tilde{\vec{x}}_n
    \tilde{\vec{x}}_n^{\top}
    \vec{u},
\]

gilt:

\[
    \frac{1}{N}
    \sum_{n=1}^{N}
    (\vec{u}^{\top}\tilde{\vec{x}}_n)^2
    =
    \vec{u}^{\top}
    \left(
        \frac{1}{N}
        \sum_{n=1}^{N}
        \tilde{\vec{x}}_n
        \tilde{\vec{x}}_n^{\top}
    \right)
    \vec{u}.
\]

Der Ausdruck in der Klammer ist die Kovarianzmatrix \(S\).
Also:

\[
    \frac{1}{N}
    \sum_{n=1}^{N}
    z_n^2
    =
    \vec{u}^{\top}S\vec{u}.
\]

\begin{definitionbox}
Die Varianz der Daten in Richtung \(\vec{u}\) ist:

\[
    \mathrm{Var}_{\vec{u}}
    =
    \vec{u}^{\top}S\vec{u},
\]

wobei

\[
    \|\vec{u}\|=1.
\]
\end{definitionbox}

\section{Erste Hauptkomponente}

Die erste Hauptkomponente (first principal component) ist die Richtung mit maximaler Varianz.

Wir suchen:

\[
    \vec{u}_1
    =
    \arg\max_{\vec{u}}
    \vec{u}^{\top}S\vec{u}
\]

unter der Nebenbedingung

\[
    \vec{u}^{\top}\vec{u}=1.
\]

\begin{definitionbox}
Die erste Hauptkomponente ist die Richtung

\[
    \vec{u}_1
\]

mit maximaler projizierter Varianz:

\[
    \vec{u}_1
    =
    \arg\max_{\|\vec{u}\|=1}
    \vec{u}^{\top}S\vec{u}.
\]
\end{definitionbox}

\section{Herleitung mit Lagrange-Multiplikatoren}

Wir maximieren:

\[
    \vec{u}^{\top}S\vec{u}
\]

unter der Nebenbedingung:

\[
    \vec{u}^{\top}\vec{u}=1.
\]

Dafür verwenden wir einen Lagrange-Multiplikator \(\lambda\).
Die Lagrange-Funktion ist:

\[
    \mathcal{L}(\vec{u},\lambda)
    =
    \vec{u}^{\top}S\vec{u}
    -
    \lambda
    (\vec{u}^{\top}\vec{u}-1).
\]

Ableiten nach \(\vec{u}\):

\[
    \nabla_{\vec{u}}\mathcal{L}
    =
    2S\vec{u}
    -
    2\lambda\vec{u}.
\]

Setze gleich \(\vec{0}\):

\[
    2S\vec{u}
    -
    2\lambda\vec{u}
    =
    \vec{0}.
\]

Teilen durch \(2\):

\[
    S\vec{u}
    =
    \lambda\vec{u}.
\]

Das ist ein Eigenwertproblem.

\begin{definitionbox}
Die Hauptkomponenten von PCA sind Eigenvektoren der Kovarianzmatrix:

\[
    S\vec{u}
    =
    \lambda\vec{u}.
\]
\end{definitionbox}

Der Wert der Varianz in Richtung \(\vec{u}\) ist:

\[
    \vec{u}^{\top}S\vec{u}.
\]

Wenn \(\vec{u}\) ein Eigenvektor mit Eigenwert \(\lambda\) ist, gilt:

\[
    S\vec{u}
    =
    \lambda\vec{u}.
\]

Dann:

\[
    \vec{u}^{\top}S\vec{u}
    =
    \vec{u}^{\top}\lambda\vec{u}
    =
    \lambda\vec{u}^{\top}\vec{u}.
\]

Da

\[
    \vec{u}^{\top}\vec{u}=1,
\]

ist:

\[
    \vec{u}^{\top}S\vec{u}
    =
    \lambda.
\]

\begin{conceptbox}
Der Eigenwert \(\lambda\) einer Hauptkomponente ist die Varianz der Daten in dieser Richtung.
\end{conceptbox}

Die erste Hauptkomponente ist daher der Eigenvektor zum größten Eigenwert.

\section{Weitere Hauptkomponenten}

Die zweite Hauptkomponente soll ebenfalls möglichst viel Varianz erklären, aber orthogonal zur ersten Hauptkomponente sein.

Allgemein:
Die \(m\)-te Hauptkomponente ist die Richtung maximaler Varianz unter der Nebenbedingung, dass sie orthogonal zu den vorherigen Hauptkomponenten ist.

\[
    \vec{u}_m^{\top}\vec{u}_j=0
    \quad
    \text{für}
    \quad
    j<m.
\]

\begin{definitionbox}
Die Hauptkomponenten sind orthonormale Eigenvektoren der Kovarianzmatrix.

Sie erfüllen:

\[
    \vec{u}_i^{\top}\vec{u}_j
    =
    \begin{cases}
        1, & i=j, \\
        0, & i\neq j.
    \end{cases}
\]
\end{definitionbox}

Da \(S\) symmetrisch ist, können die Eigenvektoren orthonormal gewählt werden.
Wir ordnen die Eigenwerte absteigend:

\[
    \lambda_1
    \geq
    \lambda_2
    \geq
    \dots
    \geq
    \lambda_D
    \geq
    0.
\]

Die zugehörigen Eigenvektoren sind:

\[
    \vec{u}_1,\vec{u}_2,\dots,\vec{u}_D.
\]

\begin{conceptbox}
PCA nimmt die Eigenvektoren mit den größten Eigenwerten.

Diese Richtungen erklären die größte Varianz der Daten.
\end{conceptbox}

\section{Projektion auf Hauptkomponenten}

Wenn wir nur die ersten \(M\) Hauptkomponenten behalten, definieren wir:

\[
    U_M
    =
    \begin{pmatrix}
        \vec{u}_1 & \vec{u}_2 & \dots & \vec{u}_M
    \end{pmatrix}
    \in
    \mathbb{R}^{D\times M}.
\]

Ein zentrierter Datenpunkt \(\tilde{\vec{x}}\) wird projiziert durch:

\[
    \vec{z}
    =
    U_M^{\top}\tilde{\vec{x}}.
\]

Dabei ist:

\[
    \vec{z}\in\mathbb{R}^M.
\]

\begin{definitionbox}
Die PCA-Projektion auf \(M\) Hauptkomponenten ist:

\[
    \vec{z}
    =
    U_M^{\top}(\vec{x}-\vec{\mu}).
\]
\end{definitionbox}

Für alle Datenpunkte in Matrixform:

\[
    Z
    =
    \tilde{X}U_M.
\]

Hier gilt:

\[
    Z\in\mathbb{R}^{N\times M}.
\]

\section{Rekonstruktion aus der reduzierten Darstellung}

Aus der reduzierten Darstellung \(\vec{z}\) kann man den Datenpunkt näherungsweise rekonstruieren.

Für zentrierte Daten:

\[
    \hat{\vec{x}}
    =
    U_M\vec{z}.
\]

Für ursprüngliche nichtzentrierte Daten addieren wir den Mittelwert zurück:

\[
    \hat{\vec{x}}
    =
    \vec{\mu}
    +
    U_M\vec{z}.
\]

Da

\[
    \vec{z}
    =
    U_M^{\top}
    (\vec{x}-\vec{\mu}),
\]

gilt:

\[
    \hat{\vec{x}}
    =
    \vec{\mu}
    +
    U_MU_M^{\top}
    (\vec{x}-\vec{\mu}).
\]

\begin{definitionbox}
Die PCA-Rekonstruktion aus \(M\) Komponenten ist:

\[
    \hat{\vec{x}}
    =
    \vec{\mu}
    +
    U_MU_M^{\top}
    (\vec{x}-\vec{\mu}).
\]
\end{definitionbox}

Wenn

\[
    M=D,
\]

und alle Hauptkomponenten verwendet werden, dann ist die Rekonstruktion exakt:

\[
    \hat{\vec{x}}=\vec{x}.
\]

Wenn

\[
    M<D,
\]

geht Information verloren.

\section{Rekonstruktionsfehler}

Der Rekonstruktionsfehler misst, wie stark ein Datenpunkt durch die niedrigdimensionale Darstellung verändert wird:

\[
    \|\vec{x}-\hat{\vec{x}}\|^2.
\]

Für alle Datenpunkte:

\[
    E_{\mathrm{rec}}
    =
    \frac{1}{N}
    \sum_{n=1}^{N}
    \|\vec{x}_n-\hat{\vec{x}}_n\|^2.
\]

\begin{definitionbox}
Der Rekonstruktionsfehler (reconstruction error) ist:

\[
    E_{\mathrm{rec}}
    =
    \frac{1}{N}
    \sum_{n=1}^{N}
    \|\vec{x}_n-\hat{\vec{x}}_n\|^2.
\]
\end{definitionbox}

Ein wichtiges Resultat:
PCA minimiert den mittleren quadratischen Rekonstruktionsfehler unter allen \(M\)-dimensionalen linearen Unterräumen.

\begin{conceptbox}
PCA kann auf zwei äquivalente Arten verstanden werden:

\begin{itemize}
    \item Finde Richtungen maximaler Varianz.
    \item Finde einen niedrigdimensionalen linearen Unterraum mit minimalem quadratischem Rekonstruktionsfehler.
\end{itemize}
\end{conceptbox}

\section{Erklärte Varianz}

Die Eigenwerte der Kovarianzmatrix geben an, wie viel Varianz jede Hauptkomponente erklärt.

Die Gesamtvarianz der Daten ist:

\[
    \sum_{j=1}^{D}
    \lambda_j.
\]

Die durch die ersten \(M\) Hauptkomponenten erklärte Varianz ist:

\[
    \sum_{j=1}^{M}
    \lambda_j.
\]

Der erklärte Varianzanteil (explained variance ratio) ist:

\[
    r_M
    =
    \frac{
        \sum_{j=1}^{M}
        \lambda_j
    }{
        \sum_{j=1}^{D}
        \lambda_j
    }.
\]

\begin{definitionbox}
Der erklärte Varianzanteil der ersten \(M\) Hauptkomponenten ist:

\[
    r_M
    =
    \frac{
        \sum_{j=1}^{M}
        \lambda_j
    }{
        \sum_{j=1}^{D}
        \lambda_j
    }.
\]
\end{definitionbox}

\begin{examplebox}
Wenn

\[
    r_M=0{,}95,
\]

dann erklären die ersten \(M\) Hauptkomponenten

\[
    95\%
\]

der Gesamtvarianz.
\end{examplebox}

\section{Wahl der Anzahl der Komponenten}

Die Anzahl \(M\) der Hauptkomponenten ist ein Hyperparameter.

Typische Kriterien:

\begin{itemize}
    \item Wähle \(M\), sodass ein bestimmter Varianzanteil erklärt wird.
    \item Verwende einen Scree Plot.
    \item Wähle \(M\) mit Validierungsdaten für eine nachgelagerte Aufgabe.
    \item Wähle \(M\) so, dass Visualisierung möglich ist, zum Beispiel \(M=2\) oder \(M=3\).
\end{itemize}

\subsection{Varianzschwelle}

Man kann \(M\) so wählen, dass:

\[
    r_M \geq 0{,}95.
\]

Das bedeutet:
Mindestens \(95\%\) der Varianz sollen erhalten bleiben.

\subsection{Scree Plot}

Ein Scree Plot zeigt die Eigenwerte

\[
    \lambda_1,\lambda_2,\dots,\lambda_D
\]

in absteigender Reihenfolge.
Oft sucht man nach einem Knickpunkt.
Nach diesem Punkt bringen zusätzliche Komponenten nur noch wenig zusätzliche Varianz.

\begin{notebox}
Ein hoher erklärter Varianzanteil bedeutet nicht automatisch gute Leistung für jede Aufgabe.
Manchmal liegen für Klassifikation wichtige Informationen in Richtungen mit kleiner Varianz.
\end{notebox}

\section{PCA-Algorithmus}

Der PCA-Algorithmus kann so zusammengefasst werden:

\begin{enumerate}
    \item Datenmatrix \(X\) aufstellen.
    \item Empirischen Mittelwert berechnen:
    \[
        \vec{\mu}
        =
        \frac{1}{N}
        \sum_{n=1}^{N}
        \vec{x}_n.
    \]
    \item Daten zentrieren:
    \[
        \tilde{\vec{x}}_n
        =
        \vec{x}_n-\vec{\mu}.
    \]
    \item Kovarianzmatrix berechnen:
    \[
        S
        =
        \frac{1}{N}
        \tilde{X}^{\top}\tilde{X}.
    \]
    \item Eigenwertproblem lösen:
    \[
        S\vec{u}_j
        =
        \lambda_j\vec{u}_j.
    \]
    \item Eigenwerte absteigend sortieren.
    \item Die ersten \(M\) Eigenvektoren wählen:
    \[
        U_M
        =
        (\vec{u}_1,\dots,\vec{u}_M).
    \]
    \item Daten projizieren:
    \[
        Z
        =
        \tilde{X}U_M.
    \]
\end{enumerate}

\begin{definitionbox}
PCA-Projektion für einen neuen Datenpunkt:

\[
    \vec{z}
    =
    U_M^{\top}(\vec{x}-\vec{\mu}).
\]

Wichtig:
Der Mittelwert \(\vec{\mu}\) und die Matrix \(U_M\) werden auf den Trainingsdaten gelernt.
\end{definitionbox}

\section{PCA und Data Leakage}

Wenn PCA als Vorverarbeitung für ein überwachtes Lernproblem verwendet wird, muss PCA nur auf den Trainingsdaten gefittet werden.

Falsch wäre:

\[
    \text{PCA auf allen Daten fitten}
    \quad
    \rightarrow
    \quad
    \text{Train/Test Split}.
\]

Richtig ist:

\[
    \text{Train/Test Split}
    \quad
    \rightarrow
    \quad
    \text{PCA auf Training fitten}
    \quad
    \rightarrow
    \quad
    \text{Training und Test transformieren}.
\]

\begin{examtipbox}
PCA muss bei einem ML-Workflow wie andere Vorverarbeitung behandelt werden.

Fit auf Trainingsdaten.
Transformation mit denselben PCA-Komponenten auf Validierungs- und Testdaten.
\end{examtipbox}

\section{Standardisierung vor PCA}

PCA ist empfindlich gegenüber Skalierung.
Wenn ein Merkmal sehr große numerische Werte hat, kann es die Hauptkomponenten dominieren.

Beispiel:
Ein Merkmal wird in Metern gemessen, ein anderes in Millimetern.
Dann ist die Varianz des Millimeter-Merkmals numerisch viel größer, obwohl es nicht unbedingt wichtiger ist.

Eine Standardisierung ist:

\[
    x_j'
    =
    \frac{x_j-\mu_j}{\sigma_j}.
\]

\begin{notebox}
Ob man vor PCA standardisieren sollte, hängt vom Problem ab.

Wenn die Skalen der Merkmale bedeutungsvoll sind, kann Zentrierung reichen.
Wenn Merkmale unterschiedliche Einheiten oder Größenordnungen haben, ist Standardisierung oft sinnvoll.
\end{notebox}

\section{PCA und Dekorrelation}

Wenn wir alle Hauptkomponenten verwenden, sind die projizierten Koordinaten unkorreliert.

Sei

\[
    U
    =
    (\vec{u}_1,\dots,\vec{u}_D)
\]

die Matrix aller Eigenvektoren.
Dann gilt:

\[
    U^{\top}SU
    =
    \Lambda,
\]

wobei

\[
    \Lambda
    =
    \mathrm{diag}(\lambda_1,\dots,\lambda_D).
\]

Das bedeutet:
Die Kovarianzmatrix im PCA-Koordinatensystem ist diagonal.

\begin{conceptbox}
PCA rotiert das Koordinatensystem so, dass die Kovarianzmatrix diagonal wird.

Die neuen Koordinaten sind unkorreliert.
\end{conceptbox}

\section{Whitening}

Whitening geht einen Schritt weiter.
Nach PCA kann man jede Komponente durch die Wurzel ihres Eigenwerts teilen.

Wenn

\[
    \vec{z}
    =
    U^{\top}(\vec{x}-\vec{\mu}),
\]

dann ist eine whitened Darstellung:

\[
    \vec{z}_{\mathrm{white}}
    =
    \Lambda^{-1/2}
    \vec{z}.
\]

Dabei ist:

\[
    \Lambda^{-1/2}
    =
    \mathrm{diag}
    \left(
        \frac{1}{\sqrt{\lambda_1}},
        \dots,
        \frac{1}{\sqrt{\lambda_D}}
    \right).
\]

\begin{definitionbox}
Whitening transformiert Daten so, dass sie Mittelwert \(0\) und Kovarianzmatrix \(I\) haben.
\end{definitionbox}

\begin{notebox}
Whitening kann nützlich sein, verstärkt aber Komponenten mit kleinen Eigenwerten.
Diese können viel Rauschen enthalten.
\end{notebox}

\section{PCA über Singular Value Decomposition}

In der Praxis berechnet man PCA häufig über Singular Value Decomposition (SVD).

Für die zentrierte Datenmatrix

\[
    \tilde{X}
    \in
    \mathbb{R}^{N\times D}
\]

lautet die SVD:

\[
    \tilde{X}
    =
    A\Sigma V^{\top}.
\]

Dabei sind:

\[
    A\in\mathbb{R}^{N\times N},
\]

\[
    \Sigma\in\mathbb{R}^{N\times D},
\]

\[
    V\in\mathbb{R}^{D\times D}.
\]

Die Kovarianzmatrix ist:

\[
    S
    =
    \frac{1}{N}
    \tilde{X}^{\top}\tilde{X}.
\]

Einsetzen der SVD:

\[
    \tilde{X}^{\top}\tilde{X}
    =
    (A\Sigma V^{\top})^{\top}
    (A\Sigma V^{\top}).
\]

Also:

\[
    \tilde{X}^{\top}\tilde{X}
    =
    V\Sigma^{\top}A^{\top}
    A\Sigma V^{\top}.
\]

Da \(A\) orthogonal ist:

\[
    A^{\top}A=I.
\]

Somit:

\[
    \tilde{X}^{\top}\tilde{X}
    =
    V\Sigma^{\top}\Sigma V^{\top}.
\]

Also:

\[
    S
    =
    V
    \left(
        \frac{1}{N}
        \Sigma^{\top}\Sigma
    \right)
    V^{\top}.
\]

\begin{conceptbox}
Die rechten Singulärvektoren der zentrierten Datenmatrix sind die PCA-Richtungen.

Die Eigenwerte der Kovarianzmatrix hängen mit den quadrierten Singulärwerten zusammen.
\end{conceptbox}

Wenn \(\sigma_j\) der \(j\)-te Singulärwert ist, dann gilt:

\[
    \lambda_j
    =
    \frac{\sigma_j^2}{N}.
\]

Je nach Konvention wird statt \(N\) auch \(N-1\) verwendet.

\section{PCA für Visualisierung}

Eine häufige Anwendung von PCA ist Visualisierung.
Dazu setzt man:

\[
    M=2
\]

oder

\[
    M=3.
\]

Dann kann man Datenpunkte in zwei oder drei Dimensionen darstellen:

\[
    \vec{z}_n
    =
    U_M^{\top}
    (\vec{x}_n-\vec{\mu}).
\]

\begin{notebox}
PCA-Visualisierung zeigt nur lineare Projektionen.

Wenn die interessante Struktur nichtlinear ist, kann PCA sie möglicherweise nicht gut sichtbar machen.
\end{notebox}

Für nichtlineare Visualisierung verwendet man oft Methoden wie \(t\)-SNE oder UMAP.
Diese sind aber schwieriger zu interpretieren und haben eigene Hyperparameter.

\section{PCA für Rauschreduktion}

Wenn Rauschen vor allem in Richtungen kleiner Varianz liegt, kann PCA zur Rauschreduktion verwendet werden.
Man behält die Hauptkomponenten mit großen Eigenwerten und verwirft die mit kleinen Eigenwerten.

\[
    \vec{x}
    \quad
    \longrightarrow
    \quad
    \vec{z}
    =
    U_M^{\top}
    (\vec{x}-\vec{\mu})
    \quad
    \longrightarrow
    \quad
    \hat{\vec{x}}
    =
    \vec{\mu}
    +
    U_M\vec{z}.
\]

\begin{conceptbox}
PCA-Rauschreduktion funktioniert gut, wenn Signal in Richtungen großer Varianz liegt und Rauschen in Richtungen kleiner Varianz.

Diese Annahme ist nicht immer erfüllt.
\end{conceptbox}

\section{PCA und Kompression}

PCA kann auch als lineare Kompression verstanden werden.
Ein Datenpunkt in \(\mathbb{R}^D\) wird durch \(M\) Koeffizienten beschrieben:

\[
    \vec{z}\in\mathbb{R}^M,
\]

wobei

\[
    M<D.
\]

Zusätzlich braucht man den Mittelwert \(\vec{\mu}\) und die Komponenten \(U_M\), um zu rekonstruieren.

\[
    \hat{\vec{x}}
    =
    \vec{\mu}
    +
    U_M\vec{z}.
\]

\begin{examplebox}
Ein Bildvektor mit

\[
    D=784
\]

Pixelwerten könnte mit

\[
    M=50
\]

PCA-Koeffizienten approximiert werden.

Das ist eine starke Reduktion der Dimension.
\end{examplebox}

\section{PCA und Klassifikation}

PCA ist unüberwacht (unsupervised).
Das bedeutet:
PCA verwendet die Labels nicht.

PCA sucht Richtungen großer Varianz in den Eingaben \(\vec{x}\), nicht Richtungen, die besonders gut Klassen trennen.

\begin{notebox}
PCA kann für Klassifikation hilfreich sein, muss es aber nicht.

Eine Richtung mit kleiner Varianz kann für die Klassentrennung sehr wichtig sein.
PCA könnte diese Richtung verwerfen.
\end{notebox}

Wenn das Ziel Klassifikation ist, sollte die Anzahl der PCA-Komponenten oft mit Validierungsdaten gewählt werden.

\section{Probabilistische Interpretation}

PCA kann auch probabilistisch interpretiert werden.
In probabilistischer PCA (probabilistic PCA) nimmt man an, dass Daten aus einem niedrigdimensionalen latenten Vektor entstehen:

\[
    \vec{z}
    \in
    \mathbb{R}^M.
\]

Ein lineares generatives Modell ist:

\[
    \vec{x}
    =
    W\vec{z}
    +
    \vec{\mu}
    +
    \vec{\epsilon}.
\]

Dabei ist:

\[
    \vec{\epsilon}
    \sim
    \mathcal{N}(\vec{0},\sigma^2 I).
\]

\begin{conceptbox}
Probabilistische PCA betrachtet PCA als generatives Modell mit latenten Variablen.

Die beobachteten Daten entstehen aus niedrigdimensionalen latenten Faktoren plus Gaußschem Rauschen.
\end{conceptbox}

Die Details der probabilistischen PCA gehen über die klassische PCA hinaus, sind aber wichtig für den Zusammenhang zu latenten Variablenmodellen und Gaussian Mixture Models.

\section{Lineare und nichtlineare Dimensionsreduktion}

PCA ist linear.
Das heißt:
Die reduzierte Darstellung entsteht durch lineare Projektion.

Viele Daten liegen aber auf nichtlinearen Mannigfaltigkeiten (manifolds).
Dann kann lineare PCA unzureichend sein.

Nichtlineare Methoden sind zum Beispiel:

\begin{itemize}
    \item Kernel PCA,
    \item Isomap,
    \item Locally Linear Embedding,
    \item \(t\)-SNE,
    \item UMAP,
    \item Autoencoder.
\end{itemize}

\begin{notebox}
Nichtlineare Methoden können komplexere Strukturen erfassen, sind aber oft schwerer zu interpretieren.

PCA bleibt wichtig, weil sie einfach, schnell, stabil und mathematisch gut verstanden ist.
\end{notebox}

\section{PCA im Vergleich zu Autoencodern}

Ein Autoencoder ist ein neuronales Netz, das lernt, Eingaben zu rekonstruieren.

Encoder:

\[
    \vec{z}
    =
    f_{\mathrm{enc}}(\vec{x}).
\]

Decoder:

\[
    \hat{\vec{x}}
    =
    f_{\mathrm{dec}}(\vec{z}).
\]

Das Training minimiert typischerweise:

\[
    \sum_{n=1}^{N}
    \|\vec{x}_n-\hat{\vec{x}}_n\|^2.
\]

Wenn Encoder und Decoder linear sind und der Fehler quadratisch ist, hängt der Autoencoder eng mit PCA zusammen.

\begin{conceptbox}
PCA kann als lineares Rekonstruktionsmodell verstanden werden.

Autoencoder verallgemeinern diese Idee mit nichtlinearen Encodern und Decodern.
\end{conceptbox}

\section{Praktische Aspekte}

\subsection{Vorverarbeitung}

Vor PCA sollte man meistens:

\begin{itemize}
    \item fehlende Werte behandeln,
    \item Daten zentrieren,
    \item je nach Problem standardisieren,
    \item Ausreißer prüfen.
\end{itemize}

\subsection{Ausreißer}

PCA kann empfindlich gegenüber Ausreißern sein.
Ein einzelner extremer Datenpunkt kann die Kovarianzmatrix und damit die Hauptkomponenten stark beeinflussen.

\subsection{Interpretierbarkeit}

Die Hauptkomponenten sind Linearkombinationen der ursprünglichen Merkmale:

\[
    z_j
    =
    \vec{u}_j^{\top}
    (\vec{x}-\vec{\mu}).
\]

Die Einträge von \(\vec{u}_j\) heißen Loadings.
Sie zeigen, welche ursprünglichen Merkmale stark zur Komponente beitragen.

\begin{definitionbox}
Loadings sind die Einträge eines Hauptkomponentenvektors \(\vec{u}_j\).

Sie beschreiben, wie stark jedes ursprüngliche Merkmal in die jeweilige Hauptkomponente eingeht.
\end{definitionbox}

\section{Häufige Fehler}

\begin{examtipbox}
Typische Fehler bei PCA und Dimensionsreduktion:

\begin{enumerate}
    \item Daten vor PCA nicht zentrieren.
    \item PCA auf allen Daten fitten und dadurch Data Leakage erzeugen.
    \item PCA-Komponenten mit ursprünglichen Merkmalen verwechseln.
    \item Denken, PCA verwende Labels.
    \item Hohe erklärte Varianz mit guter Klassifikationsleistung gleichsetzen.
    \item Standardisierung vergessen, obwohl Merkmale unterschiedliche Einheiten haben.
    \item Eigenwerte und Eigenvektoren verwechseln.
    \item Die Dimensionen von \(U_M\), \(\vec{z}\) und \(\vec{x}\) verwechseln.
    \item PCA als nichtlineare Methode interpretieren.
    \item Rekonstruktionsfehler und Klassifikationsfehler verwechseln.
    \item Scree Plot als exakte mathematische Regel statt als Diagnosewerkzeug verstehen.
\end{enumerate}
\end{examtipbox}

\section{Mini-Aufgaben}

\subsection{Aufgabe 1: Zentrierung}

Gegeben seien eindimensionale Daten:

\[
    x_1=2,
    \qquad
    x_2=4,
    \qquad
    x_3=6.
\]

Berechne den Mittelwert und die zentrierten Daten.

\begin{examplebox}
Der Mittelwert ist:

\[
    \mu
    =
    \frac{2+4+6}{3}
    =
    4.
\]

Die zentrierten Daten sind:

\[
    \tilde{x}_1=2-4=-2,
\]

\[
    \tilde{x}_2=4-4=0,
\]

\[
    \tilde{x}_3=6-4=2.
\]
\end{examplebox}

\subsection{Aufgabe 2: Erklärte Varianz}

Gegeben seien Eigenwerte:

\[
    \lambda_1=5,
    \qquad
    \lambda_2=3,
    \qquad
    \lambda_3=2.
\]

Wie viel Varianz erklären die ersten zwei Hauptkomponenten?

\begin{examplebox}
Die Gesamtvarianz ist:

\[
    \lambda_1+\lambda_2+\lambda_3
    =
    5+3+2
    =
    10.
\]

Die ersten zwei Komponenten erklären:

\[
    \lambda_1+\lambda_2
    =
    5+3
    =
    8.
\]

Der erklärte Varianzanteil ist:

\[
    r_2
    =
    \frac{8}{10}
    =
    0{,}8.
\]

Also erklären die ersten zwei Hauptkomponenten

\[
    80\%
\]

der Varianz.
\end{examplebox}

\subsection{Aufgabe 3: Projektion}

Gegeben sei

\[
    \vec{x}
    =
    \begin{pmatrix}
        3 \\
        4
    \end{pmatrix},
    \qquad
    \vec{\mu}
    =
    \begin{pmatrix}
        1 \\
        1
    \end{pmatrix},
\]

und die Hauptkomponente

\[
    \vec{u}_1
    =
    \frac{1}{\sqrt{2}}
    \begin{pmatrix}
        1 \\
        1
    \end{pmatrix}.
\]

Berechne die eindimensionale PCA-Koordinate \(z_1\).

\begin{examplebox}
Zuerst zentrieren:

\[
    \tilde{\vec{x}}
    =
    \vec{x}-\vec{\mu}
    =
    \begin{pmatrix}
        3 \\
        4
    \end{pmatrix}
    -
    \begin{pmatrix}
        1 \\
        1
    \end{pmatrix}
    =
    \begin{pmatrix}
        2 \\
        3
    \end{pmatrix}.
\]

Dann projizieren:

\[
    z_1
    =
    \vec{u}_1^{\top}\tilde{\vec{x}}
    =
    \frac{1}{\sqrt{2}}
    \begin{pmatrix}
        1 & 1
    \end{pmatrix}
    \begin{pmatrix}
        2 \\
        3
    \end{pmatrix}.
\]

Also:

\[
    z_1
    =
    \frac{1}{\sqrt{2}}
    (2+3)
    =
    \frac{5}{\sqrt{2}}.
\]
\end{examplebox}

\subsection{Aufgabe 4: Rekonstruktion}

Verwende die Werte aus Aufgabe 3 und rekonstruiere \(\hat{\vec{x}}\) aus einer Komponente.

\begin{examplebox}
Die Rekonstruktion ist:

\[
    \hat{\vec{x}}
    =
    \vec{\mu}
    +
    \vec{u}_1z_1.
\]

Einsetzen:

\[
    \vec{u}_1z_1
    =
    \frac{1}{\sqrt{2}}
    \begin{pmatrix}
        1 \\
        1
    \end{pmatrix}
    \frac{5}{\sqrt{2}}
    =
    \frac{5}{2}
    \begin{pmatrix}
        1 \\
        1
    \end{pmatrix}.
\]

Also:

\[
    \vec{u}_1z_1
    =
    \begin{pmatrix}
        2{,}5 \\
        2{,}5
    \end{pmatrix}.
\]

Damit:

\[
    \hat{\vec{x}}
    =
    \begin{pmatrix}
        1 \\
        1
    \end{pmatrix}
    +
    \begin{pmatrix}
        2{,}5 \\
        2{,}5
    \end{pmatrix}
    =
    \begin{pmatrix}
        3{,}5 \\
        3{,}5
    \end{pmatrix}.
\]
\end{examplebox}

\subsection{Aufgabe 5: Dimensionen}

Gegeben sei eine Datenmatrix mit

\[
    N=100
\]

Datenpunkten und

\[
    D=20
\]

Merkmalen.
PCA soll auf

\[
    M=3
\]

Dimensionen reduzieren.
Welche Dimensionen haben \(\tilde{X}\), \(S\), \(U_M\) und \(Z\)?

\begin{examplebox}
Die zentrierte Datenmatrix ist:

\[
    \tilde{X}\in\mathbb{R}^{100\times 20}.
\]

Die Kovarianzmatrix ist:

\[
    S\in\mathbb{R}^{20\times 20}.
\]

Die ersten \(M=3\) Hauptkomponenten bilden:

\[
    U_M\in\mathbb{R}^{20\times 3}.
\]

Die projizierten Daten sind:

\[
    Z=\tilde{X}U_M
    \in
    \mathbb{R}^{100\times 3}.
\]
\end{examplebox}

\section{Zusammenfassung}

\begin{itemize}
    \item Dimensionsreduktion (dimensionality reduction) bildet Daten aus \(\mathbb{R}^D\) nach \(\mathbb{R}^M\) ab, wobei \(M<D\).
    \item PCA (principal component analysis) ist eine lineare Methode zur Dimensionsreduktion.
    \item PCA sucht Richtungen maximaler Varianz.
    \item Die Daten müssen vor PCA zentriert werden:
    \[
        \tilde{\vec{x}}_n
        =
        \vec{x}_n-\vec{\mu}.
    \]
    \item Die empirische Kovarianzmatrix ist:
    \[
        S
        =
        \frac{1}{N}
        \tilde{X}^{\top}\tilde{X}.
    \]
    \item Die Varianz in Richtung \(\vec{u}\) ist:
    \[
        \vec{u}^{\top}S\vec{u}.
    \]
    \item Die Hauptkomponenten sind Eigenvektoren von \(S\):
    \[
        S\vec{u}_j
        =
        \lambda_j\vec{u}_j.
    \]
    \item Die Eigenwerte \(\lambda_j\) geben die erklärte Varianz der Komponenten an.
    \item Die PCA-Projektion ist:
    \[
        \vec{z}
        =
        U_M^{\top}(\vec{x}-\vec{\mu}).
    \]
    \item Die PCA-Rekonstruktion ist:
    \[
        \hat{\vec{x}}
        =
        \vec{\mu}
        +
        U_MU_M^{\top}
        (\vec{x}-\vec{\mu}).
    \]
    \item PCA minimiert den quadratischen Rekonstruktionsfehler unter allen linearen \(M\)-dimensionalen Unterräumen.
    \item Der erklärte Varianzanteil ist:
    \[
        r_M
        =
        \frac{\sum_{j=1}^{M}\lambda_j}
        {\sum_{j=1}^{D}\lambda_j}.
    \]
    \item PCA ist unüberwacht und verwendet keine Labels.
    \item PCA muss bei ML-Workflows nur auf Trainingsdaten gefittet werden, um Data Leakage zu vermeiden.
\end{itemize}

\section{Typische Prüfungsfragen}

\begin{examtipbox}
Mögliche Prüfungsfragen zu diesem Kapitel:

\begin{enumerate}
    \item Was ist Dimensionsreduktion?
    \item Warum kann hohe Dimensionalität problematisch sein?
    \item Was ist der Fluch der Dimensionalität?
    \item Was ist PCA?
    \item Warum müssen Daten vor PCA zentriert werden?
    \item Definiere die empirische Kovarianzmatrix.
    \item Zeige, dass die Kovarianzmatrix positiv semidefinit ist.
    \item Was bedeutet Varianz in einer Richtung?
    \item Leite das Eigenwertproblem der PCA mit Lagrange-Multiplikatoren her.
    \item Warum ist die erste Hauptkomponente der Eigenvektor zum größten Eigenwert?
    \item Warum sind Hauptkomponenten orthogonal?
    \item Wie projiziert man einen Datenpunkt auf die ersten \(M\) Hauptkomponenten?
    \item Wie rekonstruiert man einen Datenpunkt aus PCA-Koordinaten?
    \item Was ist Rekonstruktionsfehler?
    \item Erkläre die Äquivalenz zwischen maximaler Varianz und minimalem Rekonstruktionsfehler.
    \item Was ist erklärte Varianz?
    \item Wie wählt man die Anzahl der Komponenten?
    \item Warum kann PCA für Klassifikation ungeeignet sein?
    \item Wie hängt PCA mit SVD zusammen?
    \item Was ist Whitening?
    \item Warum kann PCA Data Leakage verursachen?
    \item Was ist der Unterschied zwischen PCA und Autoencodern?
\end{enumerate}
\end{examtipbox}